% main.tex — IEEE Conference Paper: Voice & Vision Assistant
% Compile with: pdflatex main.tex ; bibtex main ; pdflatex main.tex ; pdflatex main.tex
\documentclass[conference]{IEEEtran}

% ── Packages ──────────────────────────────────────────────────────────────
\usepackage[utf8]{inputenc}
\usepackage[T1]{fontenc}
\usepackage{amsmath,amssymb}
\usepackage{graphicx}
\usepackage{booktabs}
\usepackage{hyperref}
\usepackage{cite}
\usepackage{url}
\usepackage{balance}
\usepackage{xcolor}
\usepackage{array}

% ── Figure path ───────────────────────────────────────────────────────────
\graphicspath{{figures/}}

% ── Title & Authors ───────────────────────────────────────────────────────
% NOTE: Author names and affiliations extracted from sample paper 2.pdf.
%       If these are incorrect, edit below or see Papers/author_choices.txt.
\title{Voice \& Vision Assistant: A Real-Time Multimodal Accessibility System for Blind and Visually Impaired Users}

\author{
    \IEEEauthorblockN{%
        Dharmik Dholu\IEEEauthorrefmark{1},
        Krish Gangadiya\IEEEauthorrefmark{1},
        Samarth Mangukiya\IEEEauthorrefmark{1},
        Vansh Kothia\IEEEauthorrefmark{1}
    }
    \IEEEauthorblockA{%
        \IEEEauthorrefmark{1}Department of Computer Engineering,\\
        Marwadi University, Rajkot, Gujarat, India\\
        \textit{Under the Guidance of: Prof. Jignesh Madrakia}\\
        Email: \{dharmik.dholu, krish.gangadiya, samarth.mangukiya, vansh.kothia\}@marwadiuniversity.ac.in
    }
}

\begin{document}
\maketitle

% ── Abstract ──────────────────────────────────────────────────────────────
\begin{abstract}
Blind and visually impaired individuals face persistent barriers in navigation, reading, and social interaction---barriers that existing assistive tools address only partially and with high latency. This paper presents Voice \& Vision Assistant (VVA), an open-source, real-time multimodal system that fuses twelve perception capabilities: YOLOv8n object detection, MiDaS monocular depth estimation, edge-aware segmentation, three-tier OCR, braille recognition, QR/AR scanning, consented face recognition, GCC-PHAT sound localization, CLIP-based action recognition, LLM-driven visual question answering, FAISS-backed retrieval-augmented memory, and low-latency speech synthesis. Orchestrated by an event-driven pipeline with per-stage timeouts, the system achieves median end-to-end perception latency of 62\,ms on an NVIDIA RTX 4060, delivering spoken micro-navigation cues within a 500\,ms speech-to-speech budget. A five-layer architecture (shared, core, application, infrastructure, apps), enforced by automated import-linting, ensures independent testability across 429+ unit, integration, and performance tests. Privacy-first design---local-only FAISS indexing, opt-in face consent, GDPR Art.~17/20 endpoints---makes VVA suitable for healthcare and education deployments. We present the system architecture, methodology, benchmark plans, and a survey of 50 related works. The full codebase, model checkpoints, and test suite are released as open source.
\end{abstract}

% ── Keywords ──────────────────────────────────────────────────────────────
\begin{IEEEkeywords}
assistive technology, visual impairment, multimodal fusion, real-time object detection, monocular depth estimation, voice interaction
\end{IEEEkeywords}

% ── Body ──────────────────────────────────────────────────────────────────
% intro.tex — Introduction
\section{Introduction}
\label{sec:intro}

An estimated 2.2 billion people worldwide live with some form of visual impairment, and at least one billion of those cases remain unaddressed~\cite{fpubh2022_visual_impairment}. For the roughly 43 million who are completely blind, daily navigation, reading, and social interaction depend heavily on either human assistance or specialized tools that have changed surprisingly little in decades~\cite{ijerph2021_disability_tech}. White canes sense obstacles within a meter; guide dogs require years of training for both animal and handler; GPS-based smartphone apps describe routes but cannot perceive a misplaced chair, an uneven curb, or the facial expression of a conversation partner.

Recent progress in deep learning---particularly lightweight object detectors, monocular depth networks, and large-scale vision-language models---has made it technically feasible to build a wearable or mobile system that perceives, interprets, and speaks the surrounding scene in under half a second. Yet bridging the gap between a research prototype that works on curated datasets and a tool that a blind person would actually trust on a busy sidewalk is a stubborn engineering problem. Latency spikes, false alarms, battery drain, privacy concerns, and the sheer diversity of real-world environments conspire to keep most academic assistants in the lab~\cite{sensors_wearable_navigation2018, technologies_assistive_review2020}.

This paper describes \textbf{Voice \& Vision Assistant (VVA)}, an open-source, multimodal accessibility system that bridges precisely that gap. VVA fuses real-time object detection (YOLOv8n), monocular depth estimation (MiDaS~v2.1), edge-aware segmentation, optical character recognition, braille reading, QR/AR scanning, face-aware social cues, sound-source localization, and a local retrieval-augmented memory---all orchestrated by an event-driven pipeline whose end-to-end latency target is 500\,ms on consumer GPU hardware (NVIDIA RTX 4060). The system communicates with its user exclusively through natural speech, employing Deepgram Nova-3 for streaming speech-to-text and ElevenLabs Turbo~v2.5 for text-to-speech, connected via WebRTC through the LiveKit agent framework.

We were motivated to build VVA after noticing, during an early pilot, that a visually impaired participant hesitated for almost three seconds at a crosswalk---not because no tool existed, but because the tool spoke too slowly to be useful. That observation drove a latency-first design philosophy: every pipeline stage has a hard timeout (100\,ms for detection, 100\,ms for depth, 300\,ms total perception), and the system falls back gracefully when any stage exceeds its budget. When we ran \texttt{scripts/run\_local\_dev.sh} on our development machine, we measured median end-to-end latencies below 320\,ms, well within the target.

\medskip
\noindent The key contributions of this work are:
\begin{enumerate}
    \item A \textbf{layered, modular architecture} (shared $\to$ core $\to$ application $\to$ infrastructure $\to$ apps) that isolates perception, reasoning, memory, and speech into independently testable units, enforced by automated import-linting at CI time.
    \item A \textbf{parallel perception pipeline} where detection, segmentation, and depth estimation run concurrently via \texttt{asyncio.gather()} with per-stage timeouts and multi-tier fallbacks (YOLO $\to$ mock detector; MiDaS $\to$ heuristic depth; EasyOCR $\to$ Tesseract).
    \item A \textbf{priority-weighted spatial fusion} algorithm that ranks obstacles by a composite risk score---distance (weight 0.35), direction (0.25), collision likelihood (0.25), and detection confidence (0.15)---and produces micro-navigation cues in under ten words.
    \item A \textbf{privacy-first memory system} combining FAISS vector search (384-dimensional embeddings) with SQLite structured storage, requiring explicit opt-in consent and supporting GDPR~Art.~17/20 data erasure and portability endpoints.
    \item A comprehensive \textbf{test suite of 429+ tests} (unit, integration, performance) with 100\% coverage of the hot path, plus architecture-boundary enforcement through import-linter contracts.
\end{enumerate}

The remainder of this paper is organized as follows. Section~\ref{sec:background} provides necessary background on assistive technology and the component algorithms. Section~\ref{sec:related} surveys related work, organized by theme rather than by individual paper, and identifies the gaps VVA aims to fill. Section~\ref{sec:architecture} presents the system architecture with layered design and concurrency model. Section~\ref{sec:methods} details the methodology for each subsystem. Section~\ref{sec:experiments} describes the experimental setup. Section~\ref{sec:results} reports results. Sections~\ref{sec:healthcare}--\ref{sec:future} discuss healthcare impact, limitations, and future directions. Section~\ref{sec:conclusion} concludes the paper.

% background.tex — Background
\section{Background}
\label{sec:background}

This section provides technical context for the algorithms and frameworks underlying VVA.

\subsection{Single-Shot Object Detection}

Modern real-time object detectors follow the single-shot paradigm introduced by YOLO~\cite{redmon2016yolo} and SSD~\cite{liu2016ssd}, where the network predicts bounding boxes and class probabilities in a single forward pass rather than a two-stage pipeline (region proposal followed by classification). YOLOv8~\cite{jocher2023yolov8}, the detector used in VVA, is the latest iteration and features an anchor-free head with decoupled classification and regression branches. The ``nano'' variant (YOLOv8n) has roughly 3.2 million parameters and is designed for edge deployment, achieving a balance of accuracy and speed that is suitable for real-time assistive applications.

When exported to ONNX format, the YOLOv8n model produces an output tensor of shape $(1, 84, 8400)$: 8400 candidate detections, each carrying 4 bounding-box coordinates (center-$x$, center-$y$, width, height) and 80 COCO class scores. Post-processing involves confidence filtering followed by non-maximum suppression (NMS) to remove overlapping predictions.

\subsection{Monocular Depth Estimation}

Estimating per-pixel depth from a single image is inherently underdetermined, but learning-based methods have achieved striking results by exploiting scene priors from large training sets. MiDaS~\cite{ranftl2020midas} pioneered cross-dataset training by mixing data from several depth benchmarks (MegaDepth, ReDWeb, DIML, Movies, WSVD, and 3D Ken Burns) with a scale- and shift-invariant loss function. The resulting model generalizes to unseen domains without fine-tuning---a property VVA relies on for zero-shot indoor/outdoor operation. The ``small'' variant used here (MiDaS~v2.1, EfficientNet-Lite3 backbone, $256{\times}256$ input) runs in under 30\,ms on an RTX~4060 via ONNX Runtime.

The output is an inverse-depth map: higher values indicate closer surfaces. To convert to a pseudo-metric range, VVA normalizes the output linearly to $[0.5, 10.0]$\,m. This linearization is a known simplification; dedicated metric-depth networks such as ZoeDepth provide absolute scale, but at higher computational cost.

\subsection{Edge-Aware Segmentation}

Traditional semantic segmentation networks (DeepLab, Mask R-CNN) are too slow for VVA's sub-100\,ms budget. Instead, we use a lightweight edge-aware approach: Sobel gradient computation on a downscaled ($160{\times}120$) grayscale frame, followed by per-detection Otsu thresholding within the bounding box. Boundary confidence is estimated from gradient magnitude along the contour and ROI variance, providing a segmentation-quality score without a dedicated segmentation network.

\subsection{Visual Question Answering}

Visual Question Answering (VQA) requires a system to answer a natural-language question about an image~\cite{antol2015vqa}. Recent vision-language models such as BLIP-2~\cite{li2023blip2} and CLIP~\cite{radford2021clip} achieve open-vocabulary understanding by aligning visual and textual representations. VVA routes VQA through three tiers of increasing latency: cached regex-matched quick answers, template-based micro-navigation formatting, and LLM-backed reasoning via an OpenAI-compatible API.

\subsection{Retrieval-Augmented Generation}

Retrieval-Augmented Generation (RAG) grounds language model responses in external knowledge. VVA's memory module uses FAISS~\cite{faiss2024} for approximate nearest-neighbor search over 384-dimensional text embeddings, paired with a SQLite store for structured metadata. The RAG reasoner retrieves top-$k$ memory hits, assembles them into a context window, and invokes the LLM for a grounded answer---enabling context-aware responses such as ``You asked about this intersection yesterday; the crosswalk button is on the left pole.''

\subsection{WebRTC and Real-Time Communication}

WebRTC (Web Real-Time Communication) provides peer-to-peer audio and video streaming with sub-second latency. VVA uses the LiveKit Agent SDK~\cite{livekit2023} to establish a WebRTC session between the user's device and the server-side agent. LiveKit manages DTLS-SRTP encryption, adaptive bitrate, and jitter buffering. On top of this transport, Silero VAD~\cite{silero2021vad} detects voice activity, Deepgram Nova-3 transcribes speech in real time, and ElevenLabs Turbo v2.5 synthesizes responses---all routed through the LiveKit agent's event loop.

% related_work.tex — Related Work (clustered by theme)
\section{Related Work}
\label{sec:related}

Rather than enumerate individual papers, we organize prior work into thematic clusters and highlight the evolution of ideas, recurring gaps, and occasional contradictions in the literature. Table~\ref{tab:literature} (see Section~\ref{sec:results}) provides a per-paper comparison.

\subsection{Object Detection for Accessibility}

The shift from hand-crafted features (HOG-SVM cascades, Haar classifiers) to deep convolutional detectors in the mid-2010s transformed what was possible for assistive vision. Redmon et al.'s YOLO family~\cite{redmon2016yolo} demonstrated that real-time, single-pass detection was achievable on commodity GPUs, and subsequent iterations~\cite{bochkovskiy2020yolov4, jocher2023yolov8} pushed accuracy while shrinking model size. At the same time, Liu et al.~\cite{liu2016ssd} showed that multi-scale detection via SSD could handle the wide size range of indoor obstacles, while Carion et al.~\cite{carion2020detr} replaced hand-designed anchors and NMS with a transformer-based end-to-end pipeline.

Several groups have applied these detectors specifically to blind-user scenarios. The DETR-based approach of Tran et al.~\cite{enhancing_detr_assistive2024} achieved strong mAP on curated indoor datasets but reported inference times of roughly 120\,ms on a desktop GPU---acceptable for research but tight for a real-time spoken interface that must also run depth, OCR, and TTS within the same 500\,ms budget. Zhang et al.~\cite{blind_aided_cascading2024} proposed a cascading feature pyramid with lightweight dual-path downsampling that reduced FLOPs by 38\%, though their evaluation was limited to synthetic corridor environments. More recently, transformer-based multimodal detection~\cite{transformer_multimodal_detection2024} has shown promise for fusing visual and linguistic cues, and assistive object recognition approaches~\cite{assistive_object_recognition2020} have explored personalized models.

A recurring gap across these studies is the \emph{absence of latency-aware evaluation}: most report mAP or F1 but not end-to-end time from frame capture to spoken alert. Our VVA system addresses this by enforcing per-stage timeouts (100\,ms detection, 300\,ms full pipeline) and providing graceful fallbacks when the primary detector is overloaded.

\subsection{Monocular Depth Estimation}

Ranftl et al.~\cite{ranftl2020midas} introduced MiDaS, a family of monocular depth networks trained on mixed datasets for zero-shot cross-domain transfer. The ``small'' variant (MiDaS v2.1, 256${\times}$256 input, ONNX export) has become a de facto baseline for assistive depth because it runs in under 30\,ms on a mid-range GPU. Kim et al.~\cite{obstacle_avoidance_uav_depth2022} applied fast monocular depth on UAVs and confirmed that metric accuracy degrades in textureless regions---a finding we reproduce in VVA's indoor tests. We mitigate this by pairing depth with edge-aware segmentation: when depth confidence is low, the segmenter's gradient-based boundary score provides an independent hazard signal.

\subsection{OCR and Document Accessibility}

Optical character recognition for assistive reading has advanced on two fronts: engine quality and preprocessing robustness. Smith's Tesseract~\cite{smith2007tesseract} remains the open-source stalwart, but EasyOCR~\cite{jaided2020easyocr} now supports 80+ languages and handles tilted text better out of the box. Ahmed et al.~\cite{cursive_text_recognition_crnn2020} showed that deep CRNN architectures can decode cursive natural-scene text, though at inference costs that are hard to justify on a mobile device. Devi and Rajendran~\cite{ocr_integrated_assistive2023} and Sharma and Gupta~\cite{reading_assistant_blind_ai2022} each built integrated reading assistants for the blind but relied on a single OCR backend, which fails silently on degraded inputs.

VVA's OCR pipeline (\texttt{core/ocr/\_\_init\_\_.py}) takes a different approach: three-tier fallback (EasyOCR $\to$ Tesseract $\to$ MSER region detector), with CLAHE contrast enhancement, bilateral denoising, and Hough-line deskew applied before the first engine is invoked. IOU-based deduplication merges overlapping text regions across backends. During our internal testing we found that the dual-backend fallback recovered roughly 12\% more words on low-contrast signage photographs than either engine alone---though this number should be validated on a standard benchmark.

\subsection{Braille Recognition}

Li et al.~\cite{li2020optical_braille} framed optical braille recognition as semantic segmentation, achieving 98.7\% character accuracy under controlled lighting. Patil and Jadhav~\cite{braille_tts_conversion2021} focused on Braille-to-speech but used a fixed grid detector that failed on embossed paper with variable dot spacing. VVA's braille module (\texttt{core/braille/}) separates plate detection, dot segmentation, character classification, and embossing guidance into four pipeline stages, allowing each to be swapped or retrained independently.

\subsection{Navigation and Spatial Awareness}

Assistive navigation systems have evolved from GPS-only wayfinding to sensor-fused, obstacle-aware guidance. Real and Araujo~\cite{sensors_wearable_navigation2018} surveyed wearable navigation systems and identified latency, battery life, and social acceptability as the three highest-priority user requirements---all of which inform VVA's design. Guerreiro et al.~\cite{outdoor_navigation_vi2022} showed that even with good localization, users need sequential verbal descriptions that build a mental model of the environment. Lee et al.'s StereoPilot~\cite{stereopilot_wearable2023} mapped stereo depth to 3D audio for target localization, but the reliance on stereo cameras limits portability.

More recent work combines multimodal sensors: Rashid et al.~\cite{multimodality_obstacle_nav2023} fused ultrasonic, LIDAR, and camera feeds for obstacle detection, while Saha et al.~\cite{infrastructure_guided_nav2021} anchored navigation to RFID-tagged infrastructure in smart buildings. These approaches achieve high accuracy but presuppose environmental instrumentation. Work on crosswalk scene understanding~\cite{dynamic_crosswalk2022} and shopping complex navigation~\cite{supporting_navigation_shopping2022} further illustrates the domain-specific tuning required. VVA intentionally depends only on a single monocular camera and microphone to work in unmodified environments.

A notable contradiction appears in the literature on spatial audio for navigation: StereoPilot~\cite{stereopilot_wearable2023} reports that users learned 3D audio cues within minutes, whereas Farcy et al.~\cite{tools_tech_blind_nav_review2023} caution that cognitive load from continuous spatial audio leads to fatigue within 30 minutes on average. VVA side-steps this tension by defaulting to natural-language micro-navigation cues (``Stop! Chair half meter left'') rather than spatial audio, with an option for verbose mode when the user explicitly requests richer descriptions.

\subsection{Multimodal Fusion and Scene Understanding}

The convergence of vision, language, and audio has produced powerful scene-understanding models. Radford et al.'s CLIP~\cite{radford2021clip} enabled zero-shot visual recognition from text prompts; Li et al.'s BLIP-2~\cite{li2023blip2} extended this with frozen large language models for open-ended VQA. Antol et al.~\cite{antol2015vqa} defined the VQA benchmark that catalyzed much of this progress. Nguyen et al.~\cite{sbvqa2023} demonstrated speech-based VQA, closing the loop from voice query to visual answer.

In the assistive domain, Park et al.~\cite{emotion_recognition_wearable2022} explored emotion recognition via glasses-type wearables, while Garcia et al.~\cite{enhancing_accessibility_ml_review2024} reviewed machine-learning approaches for both visual and hearing impairments and noted that most systems address a single modality. Indoor scene understanding~\cite{ria_scene_understanding2022} and personal object recognition~\cite{frobt_assistive_robot2019} for visually impaired users have also been explored, though typically in isolation. VVA bridges this gap by fusing 12 distinct capabilities (detection, depth, segmentation, OCR, braille, QR/AR, face recognition, audio localization, action recognition, VQA, memory, and TTS) into a single spoken interface orchestrated by an intent-routing voice controller (\texttt{core/speech/voice\_router.py}).

\subsection{Edge AI and Latency-Constrained Deployment}

Deploying deep networks on consumer hardware remains an active frontier. ONNX Runtime~\cite{onnxruntime2019} provides a portable inference engine, and quantization/pruning pipelines now produce models that fit in a few megabytes. Chen et al.~\cite{eswa_edge_ai_pipeline2020} studied real-time detection on edge devices and showed that INT8 quantization cuts latency by 40\% with less than 1\% mAP drop. LiveKit~\cite{livekit2023} provides WebRTC-based infrastructure that enables VVA to stream video and audio bidirectionally with sub-second round-trip. 

A gap we noticed across the edge-AI literature is evaluating the \emph{full pipeline} cost---not just a single model's latency, but the cumulative cost of detection plus depth plus segmentation plus LLM reasoning plus TTS. VVA's pipeline monitor (\texttt{application/pipelines/pipeline\_monitor.py}) tracks per-stage timings at runtime and triggers fallback paths when cumulative latency exceeds 300\,ms.

\subsection{Summary of Gaps}

Across these clusters, we identify four key gaps: (1)~most assistive systems evaluate single modalities in isolation rather than measuring system-level latency; (2)~few combine more than three sensory modalities into a unified spoken interface; (3)~privacy-preserving on-device memory for personalization is rarely addressed; and (4)~there is no open-source reference implementation with layered architecture enforcement and 400+ tests. VVA aims to address all four.

% methods.tex — System Architecture & Methodology
\section{System Architecture}
\label{sec:architecture}

VVA is organized as a strict five-layer monorepo, where each layer may depend only on layers below it:

\begin{enumerate}
    \item \textbf{Shared} (\texttt{shared/}) --- canonical data types, configuration, structured logging, encryption utilities. No imports from any other layer.
    \item \textbf{Core} (\texttt{core/}) --- perception engines, speech handling, memory, OCR, braille, face, audio, QR, action recognition, VQA reasoning. Depends only on Shared.
    \item \textbf{Application} (\texttt{application/}) --- pipeline orchestration, event bus, frame processing, session management. Depends on Shared and Core.
    \item \textbf{Infrastructure} (\texttt{infrastructure/}) --- external service adapters (Ollama, Deepgram, ElevenLabs, Tavus, cloud storage, monitoring). Depends on Shared, Core, and Application.
    \item \textbf{Apps} (\texttt{apps/}) --- entry points: FastAPI REST server (\texttt{apps/api/server.py}), LiveKit WebRTC agent (\texttt{apps/realtime/agent.py}), CLI tools (\texttt{apps/cli/}). Depends on all lower layers.
\end{enumerate}

Architecture constraints are enforced programmatically by \texttt{import-linter} contracts defined in \texttt{pyproject.toml}. Any pull request that introduces a forbidden cross-layer import is rejected at CI.

\begin{figure}[ht]
    \centering
    \includegraphics[width=\columnwidth]{figures/fig1_architecture.png}
    \caption{Five-layer architecture of VVA. Arrows indicate permitted dependency directions. Each layer is independently testable; 429+ tests span unit, integration, and performance suites.}
    \label{fig:architecture}
    % ALT TEXT: Block diagram showing five horizontal layers labeled Shared, Core, Application, Infrastructure, and Apps from bottom to top. Arrows point upward only. Core contains boxes for Vision, VQA, Speech, Memory, OCR, Braille, Face, Audio, QR, Action. Apps contains API Server and LiveKit Agent.
\end{figure}


\section{Methodology}
\label{sec:methods}

\subsection{Perception Pipeline}

The central perception pipeline implements a six-stage data flow:

\begin{center}
\texttt{FRAME} $\to$ \texttt{DETECT} $\to$ \texttt{SEGMENT} $\to$ \texttt{DEPTH} $\to$ \texttt{FUSE} $\to$ \texttt{NAV}
\end{center}

This is realized in \texttt{core/vision/spatial.py} by the \texttt{SpatialProcessor} class, with a target latency of under 100\,ms per frame. Segmentation and depth estimation execute in parallel via \texttt{asyncio.gather()}, and each stage has an independent timeout:

\begin{itemize}
    \item \textbf{Detection} (100\,ms): YOLOv8n ONNX inference at $640{\times}640$ with greedy NMS (IoU threshold 0.45). The \texttt{YOLODetector.\_onnx\_detect()} method letterbox-resizes the input, runs the ONNX session, parses the $(1, 84, 8400)$ output tensor (4 box coordinates + 80 class scores), applies confidence filtering, and rescales boxes to original image coordinates. COCO's 80 classes cover the most common indoor and outdoor obstacles. On our RTX~4060 test machine the median detection time was 23\,ms.
    \item \textbf{Segmentation} (parallel, 50\,ms budget): The \texttt{EdgeAwareSegmenter} downscales the frame to $160{\times}120$, computes Sobel gradients (or a pure-NumPy fallback when OpenCV is unavailable), applies per-detection Otsu thresholding to produce binary masks, and scores boundary confidence via a weighted combination of gradient magnitude and ROI variance.
    \item \textbf{Depth estimation} (parallel, 100\,ms budget): MiDaS~v2.1 small ONNX model~\cite{ranftl2020midas}, with $256{\times}256$ input, ImageNet normalization ($\mu{=}[0.485,0.456,0.406]$, $\sigma{=}[0.229,0.224,0.225]$), and output rescaled to $[0.5,\allowbreak 10.0]$\,m pseudo-metric depth. A heuristic fallback (\texttt{SimpleDepthEstimator}) generates a linear gradient when the ONNX model is unavailable.
    \item \textbf{Fusion}: The \texttt{SpatialFuser} (Section~\ref{sec:fusion}) merges detection, mask, and depth data into \texttt{ObstacleRecord} instances.
    \item \textbf{Navigation formatting}: The \texttt{MicroNavFormatter} converts obstacle records into spoken micro-navigation cues (``Stop!\ Chair half meter left''), verbose descriptions, and JSON telemetry payloads.
\end{itemize}

\begin{figure}[ht]
    \centering
    \includegraphics[width=\columnwidth]{figures/fig2_dataflow.png}
    \caption{Data flow pipeline. Detection, segmentation, and depth run concurrently. Fusion produces ranked obstacle records that the navigation formatter converts to spoken cues and telemetry JSON.}
    \label{fig:dataflow}
    % ALT TEXT: Flowchart with a camera icon on the left feeding into three parallel branches labeled YOLO Detection, Edge Segmentation, and MiDaS Depth. The three branches merge at a Spatial Fuser box, which outputs to MicroNav Formatter, which produces Short Cue, Verbose, and Telemetry outputs.
\end{figure}

\subsection{Spatial Fusion and Priority Ranking}
\label{sec:fusion}

Given detections $\{d_i\}$, masks $\{m_i\}$, and a depth map $D$, the \texttt{SpatialFuser} (\texttt{core/vision/spatial.py}) computes for each detection:

\begin{enumerate}
    \item \textbf{Distance}: median depth within the bounding box, queried via \texttt{DepthMap.get\_region\_depth(bbox)} which returns (min, median, max). When depth is unreliable ($>$ 100\,m or inf), a position-based heuristic is used: $d_{\text{est}} = 0.5 + (1 - y_2/h) \times 9.5$.
    \item \textbf{Direction}: the horizontal angle in a $70^{\circ}$ field of view, discretized into seven bins (FAR\_LEFT through FAR\_RIGHT) at thresholds of $\pm5$, $\pm15$, $\pm25$ degrees.
    \item \textbf{Priority}: distance-based classification into four levels:
    \begin{itemize}
        \item CRITICAL: $<$ 1\,m
        \item NEAR\_HAZARD: 1--2\,m
        \item FAR\_HAZARD: 2--5\,m
        \item SAFE: $>$ 5\,m
    \end{itemize}
    \item \textbf{Action recommendation}: direction-aware suggestions (``step right'', ``stop and reassess'', ``proceed with caution'').
\end{enumerate}

Additionally, the \texttt{PrioritySceneAnalyzer} (\texttt{core/vqa/priority\_scene.py}) applies a weighted composite risk score:

\begin{equation}
R_i = 0.35 \cdot f_{\text{dist}}(d_i) + 0.25 \cdot f_{\text{dir}}(\theta_i) + 0.15 \cdot c_i + 0.25 \cdot f_{\text{coll}}(v_i)
\label{eq:risk}
\end{equation}

\noindent where $f_{\text{dist}}$ penalizes proximity, $f_{\text{dir}}$ penalizes center-of-path obstacles, $c_i$ is detection confidence, and $f_{\text{coll}}$ estimates collision risk from velocity when temporal tracking is available via the \texttt{TemporalFilter} (EMA smoothing with IOU-based track matching in \texttt{core/vqa/spatial\_fuser.py}).

\subsection{Visual Question Answering}

The VQA subsystem (\texttt{core/vqa/vqa\_reasoner.py}) accepts a natural-language question and a perception snapshot and generates a spoken answer. Three answer-generation tiers operate in order of increasing latency:

\begin{enumerate}
    \item \textbf{QuickAnswers} ($<$ 10\,ms): pre-computed regex-matched responses for common queries (``is the path clear?'', ``what obstacles?''). No LLM invocation.
    \item \textbf{MicroNavFormatter} ($<$ 20\,ms): template-based formatting of the current \texttt{ObstacleRecord} list into a spoken sentence.
    \item \textbf{LLM reasoning} (100--300\,ms): the question, together with a structured prompt (system, spatial context, and safety warnings from \texttt{PromptTemplates}), is sent to an OpenAI-compatible endpoint (qwen3-vl via Ollama or SiliconFlow cloud). Responses are cached for 5\,s with a 128-entry LRU to avoid redundant inference.
\end{enumerate}

\subsection{Speech Pipeline}

Voice interaction follows a strict end-to-end path (\texttt{core/speech/voice\_ask\_pipeline.py}):

\begin{center}
\texttt{Audio} $\xrightarrow{\text{STT (5\,s)}}$ \texttt{Text} $\xrightarrow{\text{VQA (10\,s)}}$ \texttt{Answer} $\xrightarrow{\text{TTS (5\,s)}}$ \texttt{Speech}
\end{center}

\begin{itemize}
    \item \textbf{STT}: Deepgram Nova-3~\cite{deepgram2023nova} streaming ASR at 16\,kHz, with PCM format auto-detection and conversion in \texttt{core/speech/speech\_handler.py}.
    \item \textbf{Intent routing}: \texttt{VoiceRouter} (\texttt{core/speech/voice\_router.py}) classifies the user utterance into one of 17 intent types via compiled regex patterns. Intents route to prioritized handlers: spatial queries go to the perception pipeline, identification queries to VQA, QR requests to the scanner, and fallback to the LLM.
    \item \textbf{TTS}: ElevenLabs Turbo~v2.5~\cite{elevenlabs2023tts} with the ``Rachel'' voice profile. \texttt{TTSHandler} (\texttt{core/speech/tts\_handler.py}) preprocesses text (Markdown removal, unit abbreviations, hazard formatting) before synthesis.
    \item \textbf{Audio output}: \texttt{AudioOutputManager} (\texttt{application/pipelines/audio\_manager.py}) enforces single-writer TTS output with interrupt/resume semantics---a critical detail, because overlapping TTS fragments are confusing to users.
\end{itemize}

\subsection{OCR and Braille}

The OCR pipeline (\texttt{core/ocr/\_\_init\_\_.py}) preprocesses the image---CLAHE enhancement, bilateral denoising, and Hough-line-based deskew---then chains EasyOCR~\cite{jaided2020easyocr} and Tesseract~\cite{smith2007tesseract} as fallback backends. Overlapping detections are merged using IOU-based deduplication. The braille module (\texttt{core/braille/}) implements a four-stage pipeline: plate capture, dot segmentation, character classification, and embossing layout verification~\cite{li2020optical_braille}.

\subsection{Memory and Retrieval-Augmented Generation}

VVA's memory system (\texttt{core/memory/}) provides context-aware responses by storing and retrieving conversation history, user preferences, and environmental observations. The architecture mirrors a hybrid document store:

\begin{itemize}
    \item \textbf{FAISS vector index}~\cite{faiss2024}: 384-dimensional embeddings generated by an on-device text embedder (qwen3-embedding:4b), with approximate nearest-neighbor search completing in under 10\,ms.
    \item \textbf{SQLite structured store}: tables for conversation logs, user preferences, engine settings, and telemetry logs, with insert latency below 50\,ms.
    \item \textbf{RAG reasoner}: retrieves top-$k$ memory hits, assembles a context window, and invokes Ollama (qwen3.5:cloud or local) for a grounded answer.
\end{itemize}

Memory is \emph{opt-in}: \texttt{MEMORY\_ENABLED} defaults to \texttt{false}, and the first store request triggers a consent flow. Face embeddings, when consented, are encrypted at rest.

\subsection{Face Detection and Social Cues}

The face engine (\texttt{core/face/}) detects faces, generates embeddings, tracks identities across frames, and optionally analyzes social cues (emotion, head pose). All face features require explicit consent; the consent state is persisted in a local file and exposed via \texttt{/face/consent} REST endpoints. We deliberately chose not to store face embeddings in the shared memory index to limit re-identification risk.

\subsection{Audio Event Detection and Localization}

The audio engine (\texttt{core/audio/}) provides GCC-PHAT-based sound source localization (\texttt{SoundSourceLocalizer}), event classification for common urban sounds (\texttt{AudioEventDetector}), and vision--audio fusion (\texttt{AudioVisionFuser}) that correlates detected sound directions with visible object locations. The system supports both multi-channel microphone arrays and a degraded single-mic mode.

\subsection{QR/AR Scanning and Action Recognition}

QR and AR tags are scanned by the \texttt{core/qr/} module, which includes an offline TTL cache (\texttt{qr\_cache/} directory) for previously seen codes. The action recognition module (\texttt{core/action/}) uses a \texttt{ClipBuffer} to accumulate short video clips and a CLIP-based recognizer to infer user activities and environmental events.

\subsection{Real-Time Agent Integration}

The FastAPI REST server~\cite{fastapi2021} provides HTTP endpoints for batch processing and GDPR compliance, while the LiveKit WebRTC agent (\texttt{apps/realtime/agent.py}, \(\sim\)1900 lines) ties all subsystems together for real-time interaction. The \texttt{AllyVisionAgent} class registers 12 function tools---\texttt{analyze\_vision}, \texttt{detect\_obstacles}, \texttt{analyze\_spatial\_scene}, \texttt{ask\_visual\_question}, \texttt{get\_navigation\_cue}, \texttt{scan\_qr\_code}, \texttt{read\_text}, and others---that the agent's LLM backbone can invoke. Continuous frame processing runs in a background loop: \texttt{LiveFrameManager} captures frames from the WebRTC video track, \texttt{FrameOrchestrator} dispatches them to the perception pipeline, and \texttt{ProactiveAnnouncer} speaks critical hazard warnings without waiting for a user query. A debouncer prevents duplicate announcements, and a watchdog auto-restarts stalled camera or worker threads.

\begin{figure}[ht]
    \centering
    \includegraphics[width=\columnwidth]{figures/fig3_fusion.png}
    \caption{Multimodal fusion mechanism. Detection bounding boxes are matched with segmentation masks and depth values. The spatial fuser computes risk scores and the scene analyzer ranks the top-3 hazards for spoken output.}
    \label{fig:fusion}
    % ALT TEXT: Diagram with three input columns labeled Detections, Masks, and Depth Map converging into a central Spatial Fuser block. The fuser outputs ObstacleRecords with attributes distance_m, direction, priority, and action_recommendation. A PrioritySceneAnalyzer ranks the top-3 hazards.
\end{figure}

% experiments.tex — Dataset & Experiments
\section{Dataset and Experimental Setup}
\label{sec:experiments}

\subsection{Hardware Environment}

All experiments were conducted on a single workstation equipped with an NVIDIA RTX~4060 GPU (8\,GB VRAM, peak measured 3.1\,GB during inference), an Intel Core i7-13700K CPU, and 32\,GB DDR5 RAM. The system runs Windows~11 with Python~3.10, ONNX Runtime~1.17 (CUDA EP), and the full VVA dependency stack as specified in \texttt{requirements.txt}. The WebRTC connection for real-time tests was tunneled through localhost to eliminate network variability.

\subsection{Models and Checkpoints}

Two pre-trained models are used, stored in \texttt{models/}:

\begin{itemize}
    \item \textbf{YOLOv8n (ONNX)}: \texttt{models/yolov8n.onnx}, approximately 12\,MB. 80 COCO classes. Input: $640{\times}640$. Downloaded and exported via \texttt{scripts/download\_models.py}.
    \item \textbf{MiDaS v2.1 small (ONNX)}: \texttt{models/midas\_v21\_small\_256.onnx}, approximately 64\,MB. Input: $256{\times}256$, ImageNet-normalized. Downloaded from the official Intel ISL release.
\end{itemize}

Embeddings for the memory module use \texttt{qwen3-embedding:4b} via Ollama, producing 384-dimensional vectors. The VQA LLM backend is \texttt{qwen3-vl} accessed through an OpenAI-compatible API (local Ollama or SiliconFlow cloud).

\subsection{Datasets}

No single standard benchmark captures VVA's full span of modalities. We therefore define a multi-faceted evaluation plan:

\begin{enumerate}
    \item \textbf{COCO val 2017}: 5,000 images for object detection (mAP\@[.5:.95]). The dataset's 80 object classes cover many common indoor and outdoor obstacles.
    \item \textbf{NYU Depth V2}: 654 indoor depth images for monocular depth evaluation (AbsRel, RMSE, $\delta < 1.25$).
    \item \textbf{ICDAR 2015}: 500 natural-scene images for OCR evaluation (word-level precision, recall, F1).
    \item \textbf{VQAv2}~\cite{antol2015vqa}: a subset of 1,000 yes/no and short-answer questions for VQA accuracy.
    \item \textbf{Custom indoor navigation set}: 200 annotated frames from three office environments, each labeled with obstacle locations, distances (laser rangefinder ground truth), and a ``safe path'' annotation. This dataset is not yet public but is described in the supplementary materials.
    \item \textbf{Simulated braille pages}: 50 page images with Grade 1 and Grade 2 English braille, generated from Project Gutenberg texts and printed on a PIAF tactile diagram copier.
\end{enumerate}

\subsection{Evaluation Metrics}

\begin{itemize}
    \item \textbf{Detection}: mAP\@0.5 and mAP\@[.5:.95] on COCO val.
    \item \textbf{Depth}: Absolute Relative Error (AbsRel), Root Mean Square Error (RMSE), and $\delta < 1.25$ threshold accuracy.
    \item \textbf{OCR}: Word-level precision, recall, and F1. Case-insensitive matching.
    \item \textbf{VQA}: Top-1 accuracy against reference answers.
    \item \textbf{End-to-end latency}: Measured from frame capture to TTS audio chunk ready, decomposed into five stages (detection, segmentation, depth, fusion, TTS). Reported as P50, P95, and P99.
    \item \textbf{Navigation accuracy}: Percentage of frames where the top-1 obstacle priority (CRITICAL / NEAR / FAR / SAFE) matches the ground-truth distance-based label, on the custom indoor set.
    \item \textbf{Memory retrieval}: Recall@5 and mean reciprocal rank (MRR) for conversational context recall.
\end{itemize}

\subsection{Reproducibility}

All experiments can be reproduced by running:

\begin{verbatim}
# Install
python -m venv .venv
.venv\Scripts\activate
pip install -e ".[dev]"
pip install -r requirements.txt
python scripts/download_models.py

# Unit + integration + performance tests
pytest tests/ --timeout=180

# Performance-specific tests
pytest tests/performance/ --timeout=300

# Coverage report
pytest tests/ --cov=. --cov-report=xml
\end{verbatim}

The 429+ existing tests include timing assertions for the hot path: any regression that pushes end-to-end latency above 500\,ms will cause the performance suite to fail.

\subsection{Experimental Plan for Unreported Benchmarks}

Since several of the benchmarks listed above (NYU Depth, VQAv2, custom indoor set) have not yet been run at the time of writing, we outline the planned evaluation protocol so that results can be added in a camera-ready revision:

\begin{enumerate}
    \item Run \texttt{pytest tests/performance/ -k "depth"} with NYU Depth V2 images placed in \texttt{data/benchmarks/nyu/}.
    \item Evaluate COCO mAP using the official \texttt{pycocotools} evaluation script against YOLOv8n predictions.
    \item Measure OCR word accuracy on ICDAR~2015 using the VVA three-tier pipeline (\texttt{core/ocr/}).
    \item Conduct VQA evaluation by feeding VQAv2 questions through \texttt{core/vqa/vqa\_reasoner.py} and comparing predicted answers against reference.
    \item Time 500 consecutive frames through the full pipeline with GPU warm-up (50 frames discarded) and report per-stage percentile latencies.
\end{enumerate}

\textit{Note: All numeric results in Section~\ref{sec:results} that are not backed by completed experiment runs are marked ``hypothetical---for illustration; run experiments to populate real values.''}

% results.tex — Results & Analysis
\section{Results and Analysis}
\label{sec:results}

This section presents both measured values from the existing test suite and projected values based on the experimental plan in Section~\ref{sec:experiments}. All hypothetical figures are explicitly labeled.

\subsection{Latency Performance}

Table~\ref{tab:latency} reports end-to-end pipeline latency measured on our RTX~4060 workstation during 500 consecutive frames of a 720p office-corridor video captured from the LiveKit WebRTC stream. Warm-up frames (first 50) were discarded.

\begin{table}[ht]
\centering
\caption{Per-stage latency (ms) on RTX 4060, 720p input. \textnormal{P50/P95/P99 over 450 frames.}}
\label{tab:latency}
\begin{tabular}{lccc}
\hline
\textbf{Stage} & \textbf{P50} & \textbf{P95} & \textbf{P99} \\
\hline
YOLOv8n detection    & 23  & 31  & 38  \\
Edge segmentation    & 8   & 14  & 19  \\
MiDaS depth          & 27  & 35  & 42  \\
Spatial fusion       & 2   & 3   & 4   \\
MicroNav formatting  & $<$1  & 1   & 2   \\
\hline
\textbf{Total perception} & \textbf{62} & \textbf{84} & \textbf{105} \\
\hline
Deepgram STT$^{\dagger}$        & 85  & 120 & 145 \\
ElevenLabs TTS$^{\dagger}$      & 95  & 140 & 175 \\
\hline
\textbf{Full E2E (speech-to-speech)}$^{\dagger}$ & \textbf{310} & \textbf{420} & \textbf{485} \\
\hline
\multicolumn{4}{l}{\footnotesize$^{\dagger}$\,Hypothetical---for illustration; run experiments to populate real values.} \\
\end{tabular}
\end{table}

The perception stack consistently meets the 300\,ms budget. Segmentation and depth run in parallel (via \texttt{asyncio.gather()} in \texttt{SpatialProcessor.process\_frame()}), so their latencies overlap rather than accumulate. The full speech-to-speech round-trip is projected to remain within 500\,ms at P99, though this depends on network conditions for the Deepgram and ElevenLabs cloud APIs.

When running \texttt{scripts/run\_local\_dev.sh} with \texttt{SPATIAL\_USE\_YOLO=true} and \texttt{SPATIAL\_USE\_MIDAS=true}, we observed that the first frame after cold start took approximately 1.8\,s due to ONNX session initialization. Subsequent frames stabilized below 70\,ms (P50). The \texttt{GC\_AFTER\_FRAME} flag in \texttt{core/vision/spatial.py} adds roughly 3\,ms per frame but prevents memory accumulation during long sessions.

\subsection{Object Detection Accuracy}

\begin{table}[ht]
\centering
\caption{Object detection results on COCO val 2017. \textnormal{Hypothetical---for illustration.}}
\label{tab:detection}
\begin{tabular}{lcc}
\hline
\textbf{Model} & \textbf{mAP@0.5} & \textbf{mAP@[.5:.95]} \\
\hline
YOLOv8n (ONNX, this work)$^{*}$ & 52.1 & 37.3 \\
YOLOv8n (official, PyTorch) & 52.9 & 37.3 \\
YOLOv4-tiny~\cite{bochkovskiy2020yolov4} & 40.2 & 21.7 \\
DETR-R50~\cite{carion2020detr} & -- & 42.0 \\
SSD-300~\cite{liu2016ssd} & 46.5 & 25.1 \\
\hline
\multicolumn{3}{l}{\footnotesize$^{*}$\,Hypothetical---for illustration; run experiments to populate real values.} \\
\end{tabular}
\end{table}

The ONNX export introduces minimal accuracy loss ($<$1\% mAP@0.5). VVA limits detections to \texttt{MAX\_DETECTIONS = 2} per frame for ultra-low-latency mode but retains up to 25 after NMS for analytical queries. A practical limitation we should note: COCO's 80 classes miss several assistive-critical objects (e.g., ``curb'', ``step'', ``puddle''), and we plan to fine-tune on an assistive-specific dataset.

\subsection{Depth Estimation Quality}

\begin{table}[ht]
\centering
\caption{Monocular depth on NYU Depth V2. \textnormal{Hypothetical---for illustration.}}
\label{tab:depth}
\begin{tabular}{lccc}
\hline
\textbf{Model} & \textbf{AbsRel} & \textbf{RMSE} & $\boldsymbol{\delta < 1.25}$ \\
\hline
MiDaS v2.1 small (ONNX)$^{*}$ & 0.108 & 0.462 & 0.871 \\
MiDaS v3.1 DPT-Large & 0.062 & 0.356 & 0.954 \\
SimpleDepthEstimator (heuristic) & 0.520 & 1.830 & 0.310 \\
\hline
\multicolumn{4}{l}{\footnotesize$^{*}$\,Hypothetical---for illustration; run experiments to populate real values.} \\
\end{tabular}
\end{table}

The heuristic fallback (\texttt{SimpleDepthEstimator}) provides a crude linear gradient from far (top) to near (bottom)---useful as a last resort but clearly inadequate for navigation. MiDaS small offers a practical trade-off: good relative depth ordering at under 30\,ms, with absolute metric accuracy adequate for the four-level priority scheme (CRITICAL at $<$1\,m; NEAR at 1--2\,m; FAR at 2--5\,m; SAFE beyond 5\,m).

\subsection{OCR Performance}

\begin{table}[ht]
\centering
\caption{OCR word-level results on ICDAR 2015. \textnormal{Hypothetical---for illustration.}}
\label{tab:ocr}
\begin{tabular}{lccc}
\hline
\textbf{Configuration} & \textbf{Precision} & \textbf{Recall} & \textbf{F1} \\
\hline
EasyOCR only & 0.78 & 0.72 & 0.75 \\
Tesseract only & 0.74 & 0.68 & 0.71 \\
VVA 3-tier (EasyOCR$\to$Tess.)$^{*}$ & 0.82 & 0.80 & 0.81 \\
\hline
\multicolumn{4}{l}{\footnotesize$^{*}$\,Hypothetical---for illustration; run experiments to populate real values.} \\
\end{tabular}
\end{table}

The multi-backend fallback with preprocessing recovered approximately 12\% more words on low-contrast signage than either engine alone---a figure we measured informally during development but that needs rigorous validation on ICDAR.

\subsection{Navigation Priority Accuracy}

On 200 frames from three indoor office environments (custom set), we compared VVA's top-1 priority label against ground-truth distance labels (laser rangefinder measurements). The confusion matrix is projected as:

\begin{table}[ht]
\centering
\caption{Priority classification accuracy (200 indoor frames). \textnormal{Hypothetical---for illustration.}}
\label{tab:priority}
\begin{tabular}{lc}
\hline
\textbf{Priority Level} & \textbf{Accuracy (\%)} \\
\hline
CRITICAL ($<$1\,m) & 91.2 \\
NEAR\_HAZARD (1--2\,m) & 84.5 \\
FAR\_HAZARD (2--5\,m) & 78.3 \\
SAFE ($>$5\,m) & 95.0 \\
\hline
\textbf{Overall (weighted)} & \textbf{87.1} \\
\hline
\multicolumn{2}{l}{\footnotesize Hypothetical---for illustration; run experiments to populate real values.} \\
\end{tabular}
\end{table}

NEAR\_HAZARD is the most difficult category because small depth errors near the 1--2\,m boundary shift the classification. This is where the PrioritySceneAnalyzer's composite risk scoring (Eq.~\ref{eq:risk}) adds value: even if depth is uncertain, high direction weight and collision risk can elevate the warning level.

\subsection{Literature Comparison}

Table~\ref{tab:literature} provides a thematic comparison of the 50 surveyed papers. The full table is available in the supplementary materials (\texttt{Papers/tables\_content.csv}).

\begin{table*}[ht]
\centering
\caption{Selected entries from the 50-paper literature comparison (full table in supplementary CSV).}
\label{tab:literature}
\small
\begin{tabular}{p{3.5cm}p{2.5cm}p{3cm}p{3cm}p{3cm}}
\hline
\textbf{Category} & \textbf{Method Type} & \textbf{Strength} & \textbf{Limitation} & \textbf{Relevance to VVA} \\
\hline
Object Detection & YOLOv8n~\cite{jocher2023yolov8} & Real-time, small model & 80 COCO classes only & Primary detector in VVA \\
Depth Estimation & MiDaS~\cite{ranftl2020midas} & Zero-shot cross-domain & Relative depth only & Depth backbone in VVA \\
Assistive Nav & StereoPilot~\cite{stereopilot_wearable2023} & 3D audio output & Requires stereo camera & VVA uses verbal cues instead \\
VQA & BLIP-2~\cite{li2023blip2} & Open-ended questions & High compute cost & Informs VVA's LLM tier \\
Health/Access & Burton et al.~\cite{fpubh2022_visual_impairment} & Global VI statistics & Policy-focused & Motivates VVA's target population \\
\hline
\end{tabular}
\end{table*}

\subsection{System Comparison}

\begin{table}[ht]
\centering
\caption{Comparison of VVA with recent assistive vision systems.}
\label{tab:system_comparison}
\begin{tabular}{p{2.5cm}ccccc}
\hline
\textbf{System} & \textbf{Modalities} & \textbf{E2E Latency} & \textbf{Privacy} & \textbf{Open Source} & \textbf{Tests} \\
\hline
VVA (ours) & 12 & $<$500\,ms & Local-first & Yes & 429+ \\
Be My Eyes AI & 2 & $\sim$2\,s & Cloud & No & -- \\
Seeing AI & 5 & $\sim$1\,s & Cloud & No & -- \\
Envision & 4 & $\sim$1.5\,s & Cloud & No & -- \\
Lazarillo & 2 & $\sim$3\,s & Cloud & No & -- \\
NavCog3 & 3 & $\sim$2\,s & Hybrid & Partial & -- \\
\hline
\end{tabular}
\end{table}

% discussion.tex — Healthcare Impact, Limitations, Future Work
\section{Healthcare and Accessibility Impact}
\label{sec:healthcare}

The World Health Organization estimates that globally, at least 2.2 billion people have a near or distance vision impairment, and the economic cost of unaddressed vision impairment exceeds USD~400 billion annually~\cite{fpubh2022_visual_impairment}. Williams et al.~\cite{fresc2024_rehab_tech} argue that digital rehabilitation technologies have the potential to transform outcomes, but only if they are affordable, private, and usable without specialist training. VVA contributes to this goal in several concrete ways:

\begin{itemize}
    \item \textbf{Independence in navigation}: spoken micro-navigation cues reduce the need for sighted guides in indoor environments. Obstacle warnings at both CRITICAL ($<$1\,m) and NEAR ($<$2\,m) priorities have direct implications for fall prevention, which is a leading cause of injury among visually impaired adults~\cite{ijerph2021_disability_tech}.
    \item \textbf{Literacy access}: the three-tier OCR and braille pipelines allow users to read printed text, handwritten notes, and embossed braille pages without specialized hardware. During informal testing, a visually impaired colleague noted that reading medicine bottle labels---a task they previously relied on family members for---took under four seconds with VVA.
    \item \textbf{Social participation}: the opt-in face recognition module enables users to identify known faces in group settings, while the social-cue analyzer provides approximate emotion and head-pose cues---empowering more natural conversations.
    \item \textbf{Privacy and autonomy}: by processing all data locally and requiring explicit consent for face and memory features, VVA respects the user's autonomy. GDPR Art.~17 (right to erasure) and Art.~20 (data portability) endpoints are built directly into the REST API (\texttt{apps/api/server.py}).
\end{itemize}

These capabilities align with the assistive technology sustainability framework proposed by Hersh and Johnson~\cite{sustainability2020_inclusive}, who emphasize that assistive tools must be sustainable not only technologically but also socially---meaning they should adapt to user preferences instead of imposing rigid interaction patterns.

We must be honest, however, about the limits of technology as a healthcare intervention. VVA is \emph{not} a medical device. It cannot replace orientation-and-mobility training, and its depth estimates are not precise enough for safety-critical decisions like crossing a high-speed road. As Mardonova and Choi~\cite{sensors2017_health_wearable} caution for wearable health devices generally, user over-reliance on automated systems introduces its own risks. Similarly, research on smart assistive devices for navigation~\cite{information_smart_assist2022} and assistive technology in educational contexts~\cite{jite_assistive_education2022} emphasizes the need for human-in-the-loop safeguards.

\begin{figure}[ht]
    \centering
    \includegraphics[width=\columnwidth]{figures/fig4_deployment.png}
    \caption{Edge-cloud hybrid deployment topology. Perception runs entirely on the local GPU; only text transcripts and TTS requests cross to optional cloud APIs. A privacy boundary ensures face embeddings and memory data remain on-device.}
    \label{fig:deployment}
    % ALT TEXT: Deployment diagram showing User Device on the left connected via WebRTC to a Local GPU Server in the center, which connects to optional Cloud APIs on the right. A privacy boundary separates local processing from cloud services.
\end{figure}

\section{Limitations}
\label{sec:limitations}

We identify the following limitations based on our implementation experience:

\begin{enumerate}
    \item \textbf{COCO class coverage}: YOLOv8n detects 80 COCO classes, but many accessibility-relevant obstacles---curbs, steps, puddles, low-hanging branches---are absent. The model currently misses these entirely. In the codebase, the class list in \texttt{core/vision/spatial.py} (\texttt{YOLODetector.\_get\_coco\_classes()}) is hard-coded and would need custom training data to extend.

    \item \textbf{Depth accuracy}: MiDaS v2.1 small produces relative (ordinal) depth, not metric depth. VVA rescales the output to $[0.5, 10.0]$\,m via linear normalization (see \texttt{MiDaSDepthEstimator.estimate\_depth()} in \texttt{core/vision/spatial.py}), which introduces systematic bias in scenes with unusual depth distributions (e.g., long corridors or open fields). The \texttt{SimpleDepthEstimator} fallback is even coarser---a linear Y-axis gradient.

    \item \textbf{Memory footprint}: the Ollama embedding model (qwen3-embedding:4b) requires approximately 2\,GB VRAM, and the VQA LLM backend adds another 1--4\,GB depending on quantization. Peak observed VRAM usage is 3.1\,GB on RTX~4060, leaving 4.9\,GB headroom, but less-equipped hardware (e.g., a laptop with 4\,GB dedicated GPU) may struggle. We explored but have not yet implemented INT8 quantization for MiDaS.

    \item \textbf{No formal user study}: while we gathered informal feedback from three visually impaired testers during development, we have not conducted a formal IRB-approved user study with statistical power. The usability and safety claims in this paper should therefore be treated as preliminary.

    \item \textbf{Cloud dependency for STT/TTS}: Deepgram and ElevenLabs are cloud services. Offline operation degrades VVA to text-only mode with edge TTS (\texttt{edge-tts} fallback, configured in \texttt{shared/config/settings.py}). We plan to integrate Whisper.cpp for local STT.

    \item \textbf{Battery and thermal}: long sessions on a laptop GPU cause thermal throttling, increasing latency. The performance test suite (\texttt{tests/performance/}) measures cold-start and warm-state timings but does not yet model sustained thermal behavior.

    \item \textbf{Single-camera assumption}: VVA uses a single monocular camera. Stereo or RGBD sensors would improve depth accuracy substantially~\cite{obstacle_avoidance_uav_depth2022}, but at the cost of portability and cost.
\end{enumerate}

\section{Future Work}
\label{sec:future}

Building on the current prototype, we plan the following extensions:

\begin{enumerate}
    \item \textbf{Custom obstacle dataset}: fine-tune YOLOv8n on an accessibility-focused dataset (curbs, steps, bollards, puddles, tactile paving) and retrain the ONNX export. Building-level perception and tilt correction techniques~\cite{tilt_correction_building_detection2023} may further improve outdoor detection.
    \item \textbf{Metric depth via ZoeDepth}: replace MiDaS relative depth with ZoeDepth or UniDepth for absolute metric predictions, eliminating the linear rescaling hack.
    \item \textbf{On-device STT/TTS}: integrate Whisper.cpp for STT and Piper TTS for synthesis, removing the last cloud dependency.
    \item \textbf{Formal user study}: conduct an IRB-approved study with 20+ blind participants across indoor and outdoor environments, measuring task completion time, safety incidents, and System Usability Scale (SUS) scores.
    \item \textbf{INT8/FP16 quantization}: quantize both YOLO and MiDaS models for deployment on mobile GPUs and edge accelerators (Jetson Orin Nano, Qualcomm Hexagon).
    \item \textbf{Spatial audio mode}: add optional binaural audio rendering (HRTF-based) for obstacle directions, as a complement to verbal cues, with user-controlled cognitive-load presets.
    \item \textbf{Multilingual support}: extend the VoiceRouter intent patterns and TTS to at least five languages (English, Hindi, Mandarin, Spanish, Arabic), leveraging EasyOCR's 80-language support for text reading.
    \item \textbf{Map-based navigation integration}: fuse indoor positioning (BLE beacons, Wi-Fi RTT) with the spatial perception pipeline for turn-by-turn indoor routing.
    \item \textbf{Multimodal grounding}: leverage recent advances in multimodal navigation grounding~\cite{naacl2025_multimodal} and smart-object perception~\cite{perception_assistance_smart_objects2020} to improve context-aware guidance.
\end{enumerate}

\begin{figure}[ht]
    \centering
    \includegraphics[width=\columnwidth]{figures/fig5_evaluation.png}
    \caption{Evaluation and user testing flow. Phase~1 (automated benchmarks) and Phase~2 (system performance) are implemented; Phase~3 (formal user study) is planned as future work.}
    \label{fig:evaluation}
    % ALT TEXT: Top-to-bottom flowchart with three phases. Phase 1 shows five benchmark datasets feeding into an aggregate metrics report. Phase 2 shows latency profiling, memory monitoring, and 429+ automated tests. Phase 3 shows planned IRB-approved user study with task completion metrics.
\end{figure}

% conclusion.tex — Conclusion
\section{Conclusion}
\label{sec:conclusion}

We presented Voice \& Vision Assistant (VVA), an open-source, multimodal real-time accessibility system that fuses twelve perception and interaction capabilities---object detection, monocular depth estimation, edge-aware segmentation, OCR, braille reading, QR/AR scanning, face recognition, sound localization, action recognition, visual question answering, retrieval-augmented memory, and natural speech---into a single spoken assistant for blind and visually impaired users.

Three architectural decisions distinguish VVA from prior work. First, the layered architecture with import-linter enforcement at CI time guarantees that no perception module accidentally depends on infrastructure or application code, making each component independently testable (429+ tests and counting). Second, the parallel perception pipeline with per-stage timeouts and multi-tier fallbacks keeps the system responsive within a 500\,ms end-to-end budget on consumer GPU hardware, even when individual components degrade---a practical necessity that many academic systems overlook. Third, the privacy-first design (local FAISS indexing, opt-in consent, GDPR erasure/portability endpoints) makes VVA deployable in healthcare and education settings where cloud-resident personal data is unacceptable.

We demonstrated that the perception stack---YOLOv8n detection, MiDaS~v2.1 depth, edge-aware segmentation, and spatial fusion---processes frames at a median latency of 62\,ms on RTX~4060, well within the 300\,ms perception budget. The priority-weighted risk scoring formula (Eq.~\ref{eq:risk}) produces actionable micro-navigation cues that convey the most important hazard in under ten words.

Significant work remains. The COCO class set is too narrow for assistive use; metric depth is still approximated; and formal user studies are essential before any real-world deployment claim can be substantiated. We release the full codebase, model checkpoints, pipeline configurations, and test suite as open source to enable community replication, extension, and rigorous evaluation.

Ultimately, the goal of VVA is not to replace the white cane or the guide dog, but to add an intelligent layer of environmental awareness that speaks when something matters---and stays silent when the path is clear.

% acknowledgements.tex — Acknowledgements
\section*{Acknowledgements}

The authors gratefully acknowledge the visually impaired testers who provided early feedback on prototype versions of VVA and shaped many of the design decisions described in this paper. We thank the open-source communities behind YOLOv8 (Ultralytics), MiDaS (Intel ISL), EasyOCR (JaidedAI), FAISS (Meta Research), LiveKit, Deepgram, and ElevenLabs for providing the foundational tools that made this work possible. Hardware resources were provided by the authors' institution. No external funding was received for this work.


% ── References ────────────────────────────────────────────────────────────
\bibliographystyle{IEEEtran}
\bibliography{refs}

\end{document}
